\textbf{Problème 1.5 (Dujella, 2026).} Existe-t-il un triplet $a_1, a_2, a_3$ d'entiers positifs distincts tel que $(a_1x+1)(a_2x+1)(a_3x+1)$ soit un carré parfait pour au moins six entiers positifs $x$ ?

\medskip

\textit{Remarque.} Un exemple dû à Petričević (2023) montre que $(x+1)(55x+1)(276x+1)$ est un carré parfait pour au moins les cinq entiers positifs suivants : $x = 4, 8, 17, 89, 4870844$. D'autres exemples connus avec au moins cinq solutions sont $(a_1,a_2,a_3) = (1,2,53)$, $(1,5,26)$, $(1,10,1407)$, $(1,11,316075)$, $(1,24,97)$, $(1,84,10333)$, $(5,24,3749)$, $(1,571,27409)$ (Petričević (2023, 2026)), $(1,8,28440881)$ (Dujella (2026)), $(3,19,7858099)$ (Sadowski (2026)). On ne sait pas s'il en existe une infinité.

Il existe certaines familles infinies avec au moins quatre solutions entières positives, notamment celles avec $(a_1,a_2) = (1,2)$ (Sadowski (2026)), $(a_1,a_2) = (1,8)$, $(a_1,a_2) = (k^2-2, k^2-1)$ pour $k \geqslant 2$, $(a_1,a_2,a_3) = (a, 3a+1, 24a^2+12a+1)$, où $6a+1 = p^2$ et $3p^2+5 = 2q^2$ (Dujella (2026)), $(a_1,a_2,a_3) = (1, k^2+2, k^8+4k^6+5k^4+k^2-1)$ (Sadowski (2026)), $(a_1,a_2,a_3) = (1, k^4-3k^2+1, k^6-7k^4+13k^2-3)$ (Andrašek (2026)).

La première famille infinie à cinq solutions a été trouvée par Christopher Sadowski (2026) : $a_1=1$, $a_2=12t+7$, $a_3=12t^2+19t+8$ avec $x=1$, $x=(21t+13-3s)/2$, $x=4t+1$, $x=(21t+13+3s)/2$, $x=(16t^2+20t+5)(144t^3+276t^2+157t+24)$, où $s^2 = 33t^2+46t+17$. À partir de $(t,s)=(4,27)$ et $(t,s)=(47,274)$, on retrouve les exemples déjà connus à cinq solutions $(1,55,276)$ et $(1,571,27409)$, tandis que le couple suivant $(t,s)=(9932,57059)$ donne $(a_1,a_2,a_3)=(1,119191,1183924204)$ avec cinq solutions $x = 9932, 18704, 39729, 189881, 222743219182025507229956$.
